\section{有效值、相位与瞬态的衔接}
有效值的物理定义来自相同电阻中的平均发热。若一段时间足以覆盖完整周期，周期电压产生的平均功率是
\[
 \overline P=\frac1{RT}\int_0^Tu^2(t)\,\mathrm dt
 =\frac{U_{\rm rms}^2}{R}.
\]
因此，正负半周的平均电压可以抵消，发热却不能抵消。只有正弦波才有峰值与有效值之比 $\sqrt2$；方波、三角波及含直流分量的波形必须重新积分。若 $u=U_{\rm dc}+u_{\rm ac}$ 且 $\overline{u_{\rm ac}}=0$，展开平方后交叉项平均为零，得
$U_{\rm rms}^2=U_{\rm dc}^2+U_{{\rm ac,rms}}^2$。

本章相量采用 $u(t)=\sqrt2\operatorname{Re}[\underline U\,e^{\mathrm i\omega t}]$，所以 $\underline U$ 的模是有效值，电感、电容阻抗分别为 $\mathrm i\omega L$ 与 $1/(\mathrm i\omega C)$。相量只描述同一角频率的周期稳态。开关刚接通时，电容电压与电感电流的连续条件往往不由稳态相量解满足，必须把自由响应加回去。
例如串联RLC中，完整电荷可写为 $q=q_{\rm forced}+q_{\rm free}$，自由部分满足
\[
 L\ddot q_{\rm free}+R\dot q_{\rm free}+q_{\rm free}/C=0.
\]
当 $R>0$ 且电路稳定时，自由响应衰减后才剩下相量所表示的稳态。若理想化为 $R=0$ 并在本征频率驱动，自由响应不衰减，外力还可产生随时间增长的共振项，此时不能直接把零阻抗写成一个有限的稳态电流。

平均功率 $P=UI\cos\varphi$ 还要求 $U,I$ 都为同频正弦量的有效值。对于非正弦周期波，最稳妥的定义始终是 $P=T^{-1}\int_0^Tu(t)i(t)\,\mathrm dt$；若逐谐波分解，各谐波的有功功率分别相加，不能给所有谐波共用一个相角。
